\chapter{Gaussian Identities Used Throughout}\label{app:gauss}

This appendix collects, with proofs, the handful of Gaussian identities on
which Chapters~\ref{ch:gaussian}--\ref{ch:bp} rely. Notation: a Gaussian
density in \emph{moment} parameters is $\Nd(a; m, v)$; in \emph{natural}
parameters it is
$\exp(-\tfrac{\lambda}{2} a^2 + \theta a + \mathrm{const})$ with
$\lambda = 1/v$, $\theta = m/v$.

\section{Quadratic forms and Gaussian normalisation}

For $J \succ 0$ (symmetric positive definite) and any $h$,
\begin{equation}
  \int_{\R^K} \exp\!\Big( -\tfrac{1}{2} a^\T J a + h^\T a \Big)\, \dd a
  = \frac{(2\pi)^{K/2}}{\sqrt{\det J}}
    \exp\!\Big( \tfrac{1}{2} h^\T J^{-1} h \Big),
  \label{eq:gauss-int}
\end{equation}
by completing the square, $-\tfrac{1}{2}a^\T J a + h^\T a =
-\tfrac{1}{2}(a - J^{-1}h)^\T J (a - J^{-1}h) + \tfrac{1}{2} h^\T J^{-1}h$,
and translating. Consequently a density
$p(a) \propto \exp(-\tfrac{1}{2} a^\T J a + h^\T a)$ is the Gaussian with
mean $m = J^{-1} h$ and covariance $J^{-1}$: \emph{reading off $(J, h)$
from a log-density is the fastest route to posterior parameters}, used in
Proposition~\ref{prop:g-posterior}.

\section{Products of Gaussians in the same variable}

\begin{equation}
  \Nd(a; m_1, v_1)\, \Nd(a; m_2, v_2)
  \;\propto\; \Nd\!\big(a;\, m, v\big),
  \qquad
  \frac{1}{v} = \frac{1}{v_1} + \frac{1}{v_2},
  \qquad
  \frac{m}{v} = \frac{m_1}{v_1} + \frac{m_2}{v_2}.
  \label{eq:gauss-product}
\end{equation}
\emph{Proof.} Add the exponents; the sum of two quadratics in $a$ is a
quadratic in $a$ whose $a^2$-coefficient is
$-\tfrac{1}{2}(1/v_1 + 1/v_2)$ and whose linear coefficient is
$m_1/v_1 + m_2/v_2$; apply \eqref{eq:gauss-int} in one dimension. In
natural parameters the rule is addition:
$(\lambda, \theta) = (\lambda_1 + \lambda_2, \theta_1 + \theta_2)$ ---
which is why the BP updates of Chapter~\ref{ch:bp} are two additions per
message. $\square$

\section{Affine propagation (marginalisation)}

If $a' = \alpha a + \eta$ with $\eta \sim \Nd(0, \sigma^2)$ independent of
$a \sim \Nd(m, v)$, then
\begin{equation}
  \int \Nd(a'; \alpha a, \sigma^2)\, \Nd(a; m, v)\, \dd a
  = \Nd\!\big(a';\, \alpha m,\, \alpha^2 v + \sigma^2\big).
  \label{eq:gauss-predict}
\end{equation}
\emph{Proof.} $a'$ is a sum of independent Gaussians, hence Gaussian; its
mean is $\alpha m$ and its variance $\alpha^2 v + \sigma^2$ by linearity
of expectation and independence. Read in the reverse direction (solving
for $a$ given a Gaussian summary of $a'$), the same computation gives
$\Nd\big(a;\, m'/\alpha,\, (v' + \sigma^2)/\alpha^2\big)$, the backward
half-step of Chapter~\ref{ch:bp}. $\square$

\section{Gaussian conditioning}

For a jointly Gaussian pair with precision blocks
$\begin{psmallmatrix} J_{11} & J_{12} \\ J_{12}^\T & J_{22}
\end{psmallmatrix}$, the conditional of block $1$ given block $2$ is
Gaussian with precision $J_{11}$ (conditioning \emph{reads off} the
precision block) and mean shifted by $-J_{11}^{-1} J_{12}\, a_2$;
marginalisation, dually, is trivial in covariance parameters. This
precision/covariance duality --- conditioning is local in $J$,
marginalising is local in $\Sigma$ --- is the algebraic engine behind the
cavity recursion \eqref{eq:bp-cavity}: integrating a neighbour out of a
tridiagonal precision produces the rank-one correction
$-\beta^2/\lambda$ on the adjacent diagonal entry, by the Schur
complement
\begin{equation}
  \begin{pmatrix} \lambda & \beta \\ \beta & J_d \end{pmatrix}
  \;\longrightarrow\;
  J_d - \frac{\beta^2}{\lambda}.
  \label{eq:gauss-schur}
\end{equation}

\section{Woodbury and the two routes to the score}

The Woodbury/Sherman--Morrison--Woodbury identity,
$(A + UCV)^{-1} = A^{-1} - A^{-1}U (C^{-1} + V A^{-1} U)^{-1} V A^{-1}$,
specialised to $\Sigt = \Deltat I + e^{-2t} \Sigz$, connects the noisy
precision $\Qt$ and the posterior precision
$J = (e^{-2t}/\Deltat) I + \Qz$:
\begin{equation}
  \Qt
  = \frac{1}{\Deltat}
    \left( I - \frac{e^{-2t}}{\Deltat}\, J^{-1} \right),
  \label{eq:gauss-woodbury}
\end{equation}
which is exactly the operator form of the score--posterior identity: substituting the
posterior mean $m = J^{-1} h = J^{-1} (e^{-t}/\Deltat) x$ into
$S = (e^{-t} m - x)/\Deltat$ gives $S = -\Qt x$ identically. The
machine-precision agreement of the two score routes
(Section~\ref{sec:g-score}) is the numerical shadow of
\eqref{eq:gauss-woodbury}.
